\section{Comparing groups: which parameter changed}\label{sec:effects}

Most applications ask not for one cohort's parameters but for what differs between two groups: a treated arm and its
control, two strains, the two sexes. This section gives the general procedure; Section~\ref{sec:val-effects} gives
the evidence for its error rates. Studies that cite it should state their own design numbers (groups, housing units such as cages, vials or plates,
onset age) and any deviation.

\paragraph{Effect models.} A reference group $R$ and a test group $T$ share $\theta$ before the onset age $t_0$. From
$t_0$ on, the test group differs by a multiplier $m$ on one parameter (Section~\ref{sec:model}). We compare
\begin{itemize}[leftmargin=1.5em,itemsep=1pt]
  \item $M_0$, no difference;
  \item the \emph{singles} $M_p$, one multiplier on $p\in\{\eta,\beta,\eps,\xc,A\}$ or on time (time rescaling), and
    where relevant $\kappa$ or an added external hazard;
  \item $M_\text{full}$, the unconstrained model: from $t_0$ on, each group has its own $(\eta,\beta,\eps)$ (with $\kappa$ at
    its fixed value, and its own $h$ where the reference fit has an external hazard). For $t_0=0$ this is two independent
    fits. Every single is a special case of $M_\text{full}$, so $M_\text{full}$ is the yardstick for asking whether one
    multiplier describes the difference adequately (condition 3 below). With $\kappa$ fixed and no external hazard it
    has two more free parameters than a single.
\end{itemize}
Pairs of multipliers ($M_{pq}$) can also be fitted. Whether a pair is identifiable depends strongly on the design
(Section~\ref{sec:val-effects}). With arms that each have their own effect size, two-parameter truths were named in only
\C{itp-ident-pair-1}--\C{itp-ident-pair-8} of replicates from one to eight arms. With eight arms spread over four dosing
ages and one effect size shared by all arms, some pairs became identifiable, but production against removal
($\eta+\beta$) never did. A pair is therefore reported only when simulations at the study's own design show that it
can be named (condition 6). For $t_0>0$ the likelihood of each individual alive at $t_0$ is conditioned on survival to $t_0$
(left truncation at the switch), and the test group's parameters change on an engine segment that starts at $t_0$
(Section~\ref{sec:likelihood}). Each model is maximised in log space with L-BFGS-B on exact gradients
\citep{byrd1995limited}, from several starts; the value at the returned point is recomputed rather than taken from
the optimiser.

\paragraph{The reference parameters.} One direction in $\log(\eta,\beta,\eps)$ is not pinned by survival data
(Section~\ref{sec:identifiability}), and an effect estimate can depend on where along it the reference sits. Three
treatments are in use, and a result is reported only if it survives the one chosen:
\begin{enumerate}[leftmargin=1.5em,itemsep=1pt]
  \item \emph{Joint refit:} $\theta$ is refitted jointly with $m$ for every model, so the comparison is between
    maximised likelihoods of both groups together.
  \item \emph{Pinned ends:} the analysis is repeated with $\theta$ at the two 95\,\% ends of the reference's
    unidentified profile; a label must survive both ends.
  \item \emph{Integration:} the effect is computed along a grid of positions on the unidentified direction and
    averaged, with the identified directions integrated by a Laplace approximation; the result then depends on the
    weighting, which must be stated.
\end{enumerate}

\paragraph{Effect sizes.} The multiplier of a single is reported with its profile-likelihood interval: the set of
$m$ whose profile log-likelihood (all other parameters re-maximised) lies within $\chi^2_{1,0.95}/2=1.92$ nats of
the maximum \citep{venzon1988method}. The profile is walked outward in $\log m$ with doubling steps, the crossing is
bisected, and the interval is widened to every evaluated point inside the cut. An interval that reaches the search box
is reported as a one-sided bound.

\paragraph{Testing and naming.} For a single $p$ we write $\Lambda_p=2[\ell(M_p)-\ell(M_0)]$. A parameter is
\emph{named} as the one that changed only if every condition below holds:
\begin{enumerate}[leftmargin=1.5em,itemsep=1pt]
  \item a difference exists: a model-free test on survival (log-rank, or a permutation test over housing units
    when individuals are clustered) rejects no difference;
  \item $M_p$ beats $M_0$ at a threshold corrected for choosing the best of $K$ singles (Šidák
    \citep{sidak1967rectangular}, then Benjamini--Hochberg across the contrasts reported together
    \citep{benjamini1995controlling});
  \item $M_p$ is not rejected against $M_\text{full}$ (likelihood-ratio test with degrees of freedom equal to the
    difference in free parameters, at level 0.05);
  \item $M_p$ is clear of the runner-up single by a fixed margin (2 nats in the studies below);
  \item the label is stable: it holds under the alternative reference treatments above, under the alternative fits
    of the reference (with and without a fit window or external hazard), and when each housing unit is left out in
    turn;
  \item the design can identify it: in simulations at the study's design, data generated by a change of $p$ are
    labelled $p$ in at least 80\,\% of replicates (recall), and a $p$ label is correct in at least 80\,\% (precision);
  \item the model fits each group, not only both together. For each group $g$ the \emph{cost} of the joint fit is
    $c_g=\ell_g(\text{own})-\ell_g(\text{joint})$: the log-likelihood that group loses at the joint fit's parameters,
    against the same group fitted alone over the parameters the joint fit moves (divided by the design effect when
    individuals are clustered). A poor fit to either group, the reference or the treated one, makes it a poor model,
    so $M_p$ fails if either cost exceeds its bar. The bar is the 95th percentile of the cost when $M_p$ is true,
    simulated at the study's design and the fitted effect size; where that simulation has not been run, a model fails
    if $2c_g$ exceeds the 99th percentile of $\chi^2$ with as many degrees of freedom as the joint fit moves.
\end{enumerate}
Condition 3 tests the two costs together: with the same free parameters, $2[\ell(M_\text{full})-\ell(M_p)]=2(c_R+c_T)$.
Condition 7 tests them one group at a time, so a model that fits the treated group well by distorting the reference,
or the reverse, cannot pass on the strength of the other group. Every model's two costs are reported beside its
likelihood, and a figure of a joint fit shows, for each group, the curve of its own fit and that of the joint fit.
When condition 6 cannot be met, results are reported at the level the design supports. The first fallback is a
\emph{confusion set}: the singles that the design cannot tell apart from the best one. The second is a \emph{group}.
The parameters fall into three groups of interest by the process they describe: \emph{production} $\{\eta, A\}$,
\emph{removal} $\{\beta, \kappa\}$, and \emph{noise and threshold} $\{\eps, \xc\}$; an added external hazard $h$ belongs
to none. Confusions are mostly within a group ($\eta$ with $A$, $\beta$ with $\kappa$), so a group can often be named when
the parameter cannot (Section~\ref{sec:val-effects}). For a pair, the question becomes which groups changed. Coarser
partitions (for example production against everything else) can be used; a study states the partition it uses and
applies condition 6 to it, with recall and precision computed for groups instead of parameters. An analytic version of condition 6 avoids most simulations: for each true single at the observed effect size, the
expected value $D_{pq}$ of $2[\ell(M_p)-\ell(M_q)]$ under $M_p$ is computed exactly from the model's predicted death
distribution. $D<1$ means the two labels are confusable, and $D\ge 9$ means they separate (validated against
recovery simulations, Section~\ref{sec:val-effects}).

\paragraph{Clustered individuals.} Individuals kept together share conditions: mice in a cage, flies in a vial,
mosquitoes in a cage, worms on a plate. These housing units, not the individuals, are then the independent units, and
likelihood-ratio statistics computed as if individuals were independent overstate the evidence. Below we write
\emph{cage} for any such unit. Two corrections are in use:
\begin{enumerate}[leftmargin=1.5em,itemsep=1pt]
  \item \emph{Whole-cage permutation.} Group labels are permuted among cages (or swapped within cage, for two sexes
    housed together), and the observed statistic is compared with its permutation distribution. This is valid under
    exchangeability of cages, at the price of coarse $p$-values when there are few cages; when cages differ in size,
    the permutations can be restricted to relabellings with similar group sizes, which makes the test conservative.
  \item \emph{Design-effect scaling.} The statistic is divided by a factor $c$ estimated from the between-cage
    variance (a cluster sandwich, or a Pearson dispersion $\phi$), and the interval cut becomes $1.92\,c$. When there
    are few cages $c$ is poorly determined, so it is estimated only from cohorts with many cages and pooled.
\end{enumerate}
Scaling by $c$ is approximate. In the one design where it was calibrated by simulation it rejected a true null in a fraction
\C{fly-null-rate} of replicates at a nominal 0.05 (Section~\ref{sec:val-effects}). We therefore recommend that the
threshold in condition 2 be calibrated by simulating the null at the study's design (cages, group sizes, measured
between-cage variance), and that design-effect scaling be used only where that simulation has been done.

This calibration has not yet been run for designs with few housing units (cages, vials, plates) per group; it is planned as an update
(Section~\ref{sec:limitations}).

\paragraph{One animal driving the fit.} Before the parameters of any group are interpreted, each cohort fit is checked
for leverage: the individual with the largest first-order influence on the reported quantities ($\eta$, $\eps$,
$\eta/\eps$, $\beta\rbar/\eta$) is dropped and the cohort refitted, and the cohort is flagged if any of these quantities
moves by more than $1.5\times$ \emph{and} the refit gains more than 1.92 nats. A flagged cohort keeps its data and its fit; its statements carry a caveat. Evidence for the check's error
rates is in Section~\ref{sec:val-leverage}.

\paragraph{Comparing across cohorts.} Log-likelihood differences grow with the number of individuals. Whenever two
cohorts of different size are compared, or a numerical agreement is judged negligible, the difference is quoted per
individual ($\Delta\ell/n$); decision thresholds within one comparison remain absolute.
